\chapter{Manual - How to Run}

This project was developed for Ubuntu/Debian-based systems.
Before running any scripts, ensure that the following dependencies
are installed:

\begin{itemize}
    \item \texttt{gcc} / \texttt{g++} (C11 and C++17 support)
    \item \texttt{make} and \texttt{pkg-config}
    \item GTK4 (\texttt{libgtk-4-dev}), required only for the graphical interface
    \item At least one audio player: \texttt{paplay} or \texttt{aplay}
    for GUI and terminal playback; \texttt{ffplay} is an additional
    fallback for the GUI
    \item MATLAB, required only for the \texttt{.m} prototypes
\end{itemize}

Most of the required dependencies can be installed using the following
command:

\begin{verbatim}
sudo apt update
sudo apt install build-essential make pkg-config libgtk-4-dev \
    pulseaudio-utils alsa-utils ffmpeg
\end{verbatim}

\section{Building the Project}

To build the exercises in Part 1, open a terminal in the project's
root directory and run the following commands:

\begin{verbatim}
cd Project
make
\end{verbatim}

After the build completes successfully, the following executables
should be available in the \texttt{Project} directory:

\begin{itemize}
    \item \texttt{wav\_app}
    \item \texttt{wav\_cmp}
    \item \texttt{wav\_effects}
    \item \texttt{wav\_hist}
    \item \texttt{wav\_quant}
\end{itemize}

\section{Running the Executables}

While in the \texttt{Project} directory, run the executables as follows:

\begin{itemize}
    \item \textbf{Histogram:}
    \begin{verbatim}
./wav_hist input.wav
    \end{verbatim}
    

    \item \textbf{Quantization:}
    \begin{verbatim}
./wav_quant input.wav output.wav --bits n
    \end{verbatim}


    \item \textbf{Comparison:}
    \begin{verbatim}
./wav_cmp input.wav input_quantized.wav
    \end{verbatim}


    \item \textbf{Audio effects:}
    \begin{verbatim}
./wav_effects
    \end{verbatim}


\end{itemize}


\section{Using the Codecs}